\begin{enumerate}[itemsep = 4mm, parsep = 3mm, topsep = 5mm, partopsep = 2mm, leftmargin = 0.5cm]

\item \textbf{Analiza wykresu Cx Cz w fukncji kąta natarcia $\alpha$} 
    \akapit Na podstawie sporządzonego wykresu Cx i Cz od $\alpha$ zauważyć można że uzyskane w ramach laboratorium charakterystyki płata są zgodne z teoretycznymi przewidywaniami. Wartość współczynnika Cx odpowiadającego za siłę oporu jest zgodnie z intuicją zawsze dodatnia. Współczynnik Cz, powiązany z siłą nośną, natomiast przecina oś OX wykresu i przyjmuje ujemne wartości dla ujemnych kątów natarcia co również jest zgodne z przewidywaniami dla standardowych ( nie samostatecznych ) profili płatów. Ponadto tak jak teoretycznie zakładano uzyskany wykres Cx jest praktycznie liniowy dla małych kątów natarcia tzn poniżej ok 20 stopni.

\item \textbf{Analiza wykresu doskonałości}
   \akapit Wykres $\frac{C_x}{C_z}$ od kąta natarcia zwany również wykresem doskonałości, również pokrywa się z oczekiwaniami, posiada max w okolicy 8 stopni kąta natarcia, co oznacza że jest najbardziej wydajny aerodynamicznie dla takiego ustroju lotu.

\item \textbf{Podsumowanie}
    \akapit Uzyskane wyniki pokrywają się z przewidywaniami teoretycznymi oraz z przedstawionymi na laboratoriach historycznymi pomiarami innych profili płatów. Taka powtarzalność oraz przewidywalność wyników, nawet pomimo różnych metod obróbki danych - kiedyś ręcznego wyrównywania wag szalkowych, a dzisiaj cyfrowej analizy danych, świadczy o dużej wartości metody pomiarów z użyciem wagi aerodynamicznej nawet w epoce oprogramowania CFD.
\end{enumerate}